% Metadados do trabalho - segue convencoes do template abntex2ime

\instituicao{Instituto Militar de Engenharia}
\programapgdepartamento{Engenharia de Computa\c{c}\~ao, Engenharia Eletr\^onica e Engenharia Cartogr\'afica}
\nivelestudo{Graduação}

\titulo{Caix\~ao de Areia com Realidade Aumentada: Sistema de Visualiza\c{c}\~ao Interativa de Modelos Digitais de Eleva\c{c}\~ao para Apoio ao Treinamento Militar}

\palavraschave{realidade aumentada, modelo digital de eleva\c{c}\~ao, kinect, calibra\c{c}\~ao geom\'etrica, RANSAC, treinamento militar}
\keywords{augmented reality, digital elevation model, kinect, geometric calibration, RANSAC, military training}

% Autores: nomes e sobrenomes separados por virgula
\autores{Rafael,Raquel,Mateus,Est\^{e}v\~{a}o}{de Figueiredo Schuinki,Belchior Fa\c{c}anha Nogueira,da Silva Maldonado,Johnatas Ribeiro Batista}

% Orientadores: nomes, sobrenomes e titulos
\orientadores{Daniel,Leonardo,Bruno,Gabriel}{Rodrigues dos Santos,Assump\c{c}\~ao Moreira,Eduardo Madeira,da Cruz Fontenelle}{D.Sc.,D.Sc.,D.Sc.,M.Sc.}

\local{Rio de Janeiro}
\data{2026}
\datadefesa{08 de outubro de 2026}

\bancadeexaminadores{
    Prof. Daniel Rodrigues dos Santos - D.Sc. do IME - Presidente,
    Prof. Leonardo Assump\c{c}\~ao Moreira - D.Sc. do IME,
    Prof. Bruno Eduardo Madeira - D.Sc. do IME,
    Prof. Gabriel da Cruz Fontenelle - M.Sc. do IME
}

\preambulo{%
Projeto Final de Curso apresentado ao Curso de Gradua\c{c}\~ao em Engenharia de Computa\c{c}\~ao, Engenharia Eletr\^onica e Engenharia Cartogr\'afica do Instituto Militar de Engenharia, como requisito parcial para a obten\c{c}\~ao do t\'\i tulo de Bacharel em Engenharia de Computa\c{c}\~ao, Engenharia Eletr\^onica e Engenharia Cartogr\'afica.}
